\documentclass{article}

\textwidth=6.5in
\textheight=8.5in
\voffset=-54pt
\oddsidemargin=0in
\evensidemargin=0in
\setlength{\parskip}{8pt}
\setlength{\parindent}{0pt}

\usepackage{workbook}

\usepackage[sols]{handout}

\begin{document}

\begin{center}
  \large\textsc{Problem Set 24 - The Fundamental Theorem of Calculus, Part 1}
  
      \pspace{12pt}
  \begin{problemonly}\large Name: \rule{5cm}{0.15mm}\\ \end{problemonly}
\end{center}

For help with these problems, you may wish to read pp. 757-758 in Gottlieb.

\begin{enumerate}[topsep=0pt]
\item

\begin{grading}
10 points: 2/2/3/3
\end{grading}

  The rate that water is entering/leaving a tank is given by $f(t)=60-3t$
  gallons/minute where $t$ is measured in minutes past noon.
  Let $A(x)=\displaystyle\int_0^x f(t)\,dt$. 

  \begin{enumerate}[topsep=0pt]

  \item
    \begin{problem}[\vfill]
      Interpret $A (x)$ in words in the context of this problem, for $x\ge 0$.
    \end{problem}

    \begin{solution} % Jason Bland
      $\displaystyle A(x) = \int_0^x f(t)\,dt$ is the net change in the
      amount of water in the tank between noon and $x$ minutes after
      noon.  In other words, it is (amount of water in the tank $x$
      minutes after noon) $-$ (amount of water in the tank at noon).
    \end{solution}

  \item
    \begin{problem}[\vfill]
      For what values of $x$ is $A(x)$ increasing?
    \end{problem}

    \begin{solution} % Jason Bland
      $A(x)$ is increasing when $A'(x) = f(x) = 60 - 3x$ is positive,
      which happens for $\boxed{x < 20}$.
    \end{solution}

  \item
    \begin{problem}[\vfill]
      At what time is the water level in the tank the same as it was at
      noon? At this time, what is the value of $A(x)$?
    \end{problem}

    \begin{solution} 
      We are looking for a time $x$ when the net change in amount of water
      from noon to $x$ is 0; that is, we want $\displaystyle \int_0^x
      f(t)\,dt = 0$.  Graphically, we want to find the value of $x$ so
      that the signed area under $f(t)$ from $t = 0$ to $t = x$ is 0: 
      \begin{center}
        \begin{tikzpicture}
          \begin{axis}[ytick={60}, width=3in, height=2in, xtick={20, 40},
            xticklabels={20, $x$}]
            \addplot+[domain=0:20, plusarea] {60-3*x} \closedcycle;
            \addplot+[domain=20:40, minusarea] {60-3*x} \closedcycle;
            \addplot+[domain=-10:50, black] {60-3*x};
          \end{axis}
        \end{tikzpicture}
      \end{center}
      By symmetry, we can see that $x = 40$, which corresponds to
      \fbox{12:40 pm}.
    \end{solution}

  \item
    \begin{problem}[\vfill\newpage]
      If the water tank has $30,000$ gallons of water at noon, how much
      does it have at 1 pm?  (Be careful with units!)
    \end{problem}

    \begin{solution}
      The net change in the amount of water from noon ($t = 0$) to 1 pm
      ($t = 60$) is $\displaystyle \int_{0}^{60} f(t)\,dt$.  We can
      picture this graphically as the following signed area:
      \begin{center}
        \begin{tikzpicture}
          \begin{axis}[ytick={60}, width=3in, height=2in, xtick={0,20,60}]
            \addplot+[domain=0:20, plusarea] {60-3*x} \closedcycle;
            \addplot+[domain=20:60, minusarea] {60-3*x} \closedcycle;
            \addplot+[domain=-10:80, black] {60-3*x};
          \end{axis}
        \end{tikzpicture}
      \end{center}
      Since $f(0) = 60$ and $f(60) = -120$, we can find the area of
      the green triangle, $\frac{1}{2}(20)(60)$, and subtract the area
      of the red triangle, $\frac{1}{2}(40)(120)$, to get
      $\displaystyle \int_{0}^{60} f(t)\,dt = -1800$.
      Therefore, the net change in the amount of water from noon
      to 1 pm is $-1800$ gallons.  Since there were $30000$ gallons of
      water at noon, the amount of water at 1 pm was $30000 - 1800 =
      \boxed{28200 \text{ gallons}}$. 
    \end{solution}

  \end{enumerate}

\item 

\begin{grading}
20 points: 1/2/2/3/3/5/4/0. Part (h) should not be graded for any credit.
\end{grading}

  Let $f$ be the quadratic function whose graph is shown below.
  Let $A(x)=\displaystyle \int_2^x f(t)\,dt$.

  \begin{tikzpicture}
    \begin{axis}[xtick={-1, 1, 3}, ytick={-4}, width=2.25in]
      \addplot+[domain=-2.1:4.1]{(x-1)^2-4};
    \end{axis}
  \end{tikzpicture}

  \begin{enumerate}
  \item
    \begin{problem}[\vfill]
      Evaluate $A(2)$. 
    \end{problem}

    \begin{solution} % Jason Bland
      $\displaystyle A(2) = \int_2^2 f(t) \, dt = 0$.
    \end{solution}

  \item
    \begin{problem}[\vfill]
      Is $A(3)$ positive, negative, or zero?  Explain how you know. 
    \end{problem}

    \begin{solution} 
      $A(3) = \displaystyle \int_2^3 f(t)\,dt$, which we can visualize
      as the signed area under $f(t)$ from $t = 2$ to $t = 3$.  This
      signed area is \fbox{negative} since $f < 0$ on $(2, 3)$.
    \end{solution}

  \item
    \begin{problem}[\vfill]
      Is $A(1)$ positive, negative, or zero?  Explain how you know.
    \end{problem}

    \begin{solution}
      $A(1) = \displaystyle \int_2^1 f(t)\,dt = -\int_1^2 f(t)\,dt$.
      The integral $\displaystyle \int_1^2 f(t)\,dt$ can be visualized as
      the signed area under $f(t)$ from $t = 1$ to $t = 2$, and this
      signed area is negative.  Therefore, $A(1)$ is \fbox{positive}. 
    \end{solution}

  \item
    \begin{problem}[\vfill]
      For what values of $x$ is $A(x)$ increasing?  
    \end{problem}

    \begin{solution} % Jason Bland
      $A$ is increasing when $A'$ is positive; by FTOC 1, $A' =
      f$, and this is positive on \fbox{$(-\infty, -1)$ and $(3,
        \infty)$}.
    \end{solution}

  \item
    \begin{problem}[\vfill]
      For what values of $x$ is $A(x)$ concave up?
    \end{problem}

    \begin{solution} % Jason Bland
      $A$ is concave up when $A''$ is positive; by FTOC 1, $A'' =
      f'$, which is positive when $f$ is increasing. This happens on
      \fbox{$(1, \infty)$}.
    \end{solution}

  \item
    \begin{problem}[\nstretch{2}]
      Sketch a rough graph of $A(x)$ on the interval $[-2, 4]$. 
      \emph{Make sure it agrees with parts (a)-(e)!}
    \end{problem}

    \begin{solution} % Jason Bland
      We've seen that $A$ has an $x$-intercept at $x = 2$, and it's
      increasing on $(-\infty, -1)$, decreasing on $(-1, 3)$, and
      increasing on $(3, \infty)$.  Also, it's concave down on $(-\infty,
      1)$ and concave up on $(1, \infty)$.  Thus, its graph looks like
      this:     
      \begin{center}
        \begin{tikzpicture}
          \begin{axis}[xtick={-1, 1, 2, 3}, ytick={0}, width=2.25in, ymax=22]
            \addplot+[domain=-2.1:4.1]{(1/3) * (x - 1)^3 - 4 * (x - 1) +
              (11/3)};
          \end{axis}
        \end{tikzpicture}
      \end{center}
    \end{solution}

  \item
    \begin{problem}[\stretch{2}]

    Let $B(x)=\displaystyle\int_3^x f(t)\,dt$. 
    Sketch a rough graph of $B(x)$.
    (\emph{Hint:} What do you
      know about how $B(x)$ relates to 
      $A(x)$?  What is $B(3)$?
      You may want to look back at
      {{{Problem Set 23}, {{\#6}}}}.)
    \end{problem}

    \begin{solution} % Jason Bland
      As we saw in {{{Problem Set 23}, {{\#6}}}},
        $B(x)$ and $A(x)$ differ by a constant.
        Graphically, that means the graph of $B(x)$
        is a vertical
        shift of the graph of $A(x)$.  Since $B(3) =
        \displaystyle \int_3^3 f(t)\,dt = 0$, we must shift the graph of
        $A$ vertically until it crosses the $x$-axis at $x = 3$ to
        get the graph of $A$: 
    \begin{center}
      \begin{tikzpicture}
        \begin{axis}[xtick={-1, 1, 2, 3}, ytick={0}, width=2.25in,
          ymin=-2, ymax=22, clip=false]
          \addplot+[domain=-2.1:4.1]{(1/3) * (x - 1)^3 - 4 * (x - 1) +
            (11/3)} node[right]{$A$};
          \addplot+[domain=-2.1:4.1]{(1/3) * (x - 1)^3 - 4 * (x - 1) +
            (16/3)} node[above right]{$B$};
        \end{axis}
      \end{tikzpicture}
    \end{center}
    \end{solution}

  \item
    \begin{problem}
      Check your answers using the Area Function applet at
      \url{https://www.geogebra.org/m/cjdwbanv}.  The formula for
      $f$ is $f(t) = (t - 1)^2 - 4$, which is already entered in the
      applet.

      To see $A$, just drag the red dot labeled $x$ around; the purple
      graph on the right is $A$.

      To see $B(x)$, first change $a$ in the
      applet to be $3$, and then drag $x$ around again.
    \end{problem}

  % \item
  %   \begin{problem}
  %     Suppose that $f(t)$ represents the rate, in students per minute, at
  %     which students are entering or leaving Annenberg $t$ minutes after 8
  %     am one morning.  If there are 10 students in Annenberg at 8:03 am,
  %     sketch the graph of $S(t)$, the number of students in Annenberg $t$
  %     minutes after 8 am. 
  %   \end{problem}

  %   \begin{solution}
  %     At $t = 3$, there are $10$ students in Annenberg.  Therefore, the
  %     number of students in Annenberg $t$ minutes after 8 am is
  %     $\displaystyle 10 + \int_3^t f(u)\,du$.  That is, $S(t) = 10 + {}_3
  %     A(t)$.  Therefore, to graph $S(t)$, we just shift the graph of
  %     ${}_3 A(t)$ up by 10 units:
  %     \begin{center}
  %       \begin{tikzpicture}
  %         \begin{axis}[xtick={-1, 1, 2, 3}, ytick=\empty, width=2.25in,
  %           ymin=-2, ymax=22, clip=false]
  %           \addplot+[domain=-2.1:4.1]{(1/3) * (x - 1)^3 - 4 * (x - 1) +
  %             (11/3)} node[right]{$A$};
  %           \addplot+[domain=-2.1:4.1]{(1/3) * (x - 1)^3 - 4 * (x - 1) +
  %             (16/3)} node[above right]{$B$};
  %           \addplot+[domain=-2.1:4.1] {x^3/3 - x^2 - 3*x + 19} node[right]{$S$};
  %           \addplot+[closed, red] coordinates { (3, 10) }
  %           node[above] { $(3, 10)$ };
  %         \end{axis}
  %       \end{tikzpicture}
  %     \end{center}
  %   \end{solution}

  \end{enumerate}

\item 

\begin{grading}
7 points: 1/2/4. For part (c), they need to explain why they're able to eliminate the wrong answers, and not just why the correct answer works!
\end{grading}

    Laura is trying to find a formula for the area function 
    $A(x)=\displaystyle\int_1^x f(t)\,dt$,
    where $f(t) = \dfrac{2t}{t^2 + 1}$.

  \begin{enumerate}[topsep=0pt]
  \item
    \begin{problem}[\vfill]
      What is $A(1)$? That is, what is 
      $\displaystyle\int_1^1 f(t)\,dt$? 
    \end{problem}

    \begin{solution}
      $\displaystyle A(1) = \int_1^1 f(t)\,dt = \boxed{0}$. 
    \end{solution}

  \item
    \begin{problem}[\vfill]
      What is $A'(x)$? That is, what is
      $\displaystyle \frac{d}{dx}\left[\int_1^x f(t)\,dt\right]$? Your answer should be a function of $x$!
    \end{problem}

    \begin{solution}
      By FTOC 1, $A'(x) = f(x) = \dfrac{2x}{x^2 + 1}$. 
    \end{solution}

  \item
    \begin{problem}[\vfill\newpage]
      One of the following is a formula for $A(x)$; which one?
      (\emph{Note:} You are \textbf{not} expected to be able to directly
      find a formula for $A(x)$; however, given what you found in the
      previous parts of this problem, you should be able to eliminate all
      but one of the choices below.)
      
        \begin{enumerate}[label=\protect{\maybecircle{\Alph*.}}]
        \plainitem $\ln (x^2 + 1)$ $\vphantom{\dfrac{1}{1}}$
        \circleitem $\ln (x^2 + 1) - \ln 2$ $\vphantom{\dfrac{1}{1}}$
        \plainitem $\dfrac{2 (1 - x^2)}{(x^2 + 1)^2}$
        \plainitem $2x \arctan x$ $\vphantom{\dfrac{1}{1}}$
        \end{enumerate}
      
      Explain your reasoning.
    \end{problem}

    \begin{solution}
      We know two things about $A(x)$: its derivative is
      $\dfrac{2x}{x^2 + 1}$, and it has an $x$-intercept at $x = 1$.  Only
      (A) and (B) have derivative $\dfrac{2x}{x^2 + 1}$; of those two,
      only \fbox{(B)} has an $x$-intercept at $x = 1$.
    \end{solution}

  \end{enumerate}

\item
% 23.2.2

\begin{grading}
8 points: 2 for each part.
\end{grading}

  Here is the graph of a function $f(t)$. 
  \begin{center}
    \begin{tikzpicture}[scale=0.8]
      \begin{axis}[ xscale=0.8, xmin=-3, xmax=3, xtick=\empty, yscale=0.4,
        ymin=-0.9, ymax=0.9, ytick=\empty]
        \addplot+[no marks,black, thick,domain=-3:3,unbounded coords=jump,samples=101]{2*x/(1+x^2)^2};
      \end{axis}
    \end{tikzpicture}
  \end{center}

  \begin{enumerate}
  \item
    \begin{problem}[\vfill]
      One of the following graphs is the graph of
      $\displaystyle F(x) = \int_0^x f(t) \,dt$; which one? Explain.
    \end{problem}

    \begin{solution} % Jason Bland
      $F'(x) = f(x)$ by
        FTOC 1.  In particular,
      since $f$ is negative on $(-\infty, 0)$ and positive on $(0,
      \infty)$, $F$ is decreasing on $(-\infty, 0)$ and increasing on $(0,
      \infty)$.  This only matches (D) and (E).

      Since $\displaystyle F(0) = \int_0^0 f(t)\,dt = 0$, the graph of $F$
      must be \fbox{(D).}
    \end{solution}

  \item
    \begin{problem}[\vfill]
      One of the following graphs is the graph of $\displaystyle G(x) =
      \int_2^x f(t) \,dt$; which one? Explain.
    \end{problem}

    \begin{solution} % Jason Bland
      $G(x)$ differs
        from $F(x)$ by a constant (by
        {{{Problem Set 23}, {{\#6}}}}).  The only vertical shift
      of $F$ is \fbox{(E).}

      (Also, notice that $G(2) = \displaystyle \int_2^2 f(t)\,dt = 0$.)
    \end{solution}

  \item
    \begin{problem}[\vfill]
      One of the following graphs is the graph of $\displaystyle H(x) =
      \int_x^2 f(t)\,dt$; which one?  Explain.
    \end{problem}

    \begin{solution} % Jason Bland
      We can write $H(x) = \displaystyle \int_x^2 f(t)\,dt = - \int_2^x
      f(t)\,dt = - G(x)$.  So, to get the graph of $H$, we can flip the
      graph of $G$ over the $x$-axis, which gives us graph \fbox{(F).} 
    \end{solution}

  \item
    \begin{problem}[\vfill]
      One of the following graphs is the graph of $\displaystyle K(x) =
      \int_2^3 f(t)\,dt$; which one?  Explain.
    \end{problem}

    \begin{solution} % Jason Bland
      $K(x)$ is a positive constant, so its graph is \fbox{(J).}
    \end{solution}

  \end{enumerate}
\end{enumerate}

  Choices:
  
    \pgfplotsset{
      cycle list={blue},
      every axis/.style={
        xmin=-3.5,
        xmax=3.5,
        xtick=\empty,
        ymin=-0.9,
        ymax=0.9,
        ytick=\empty,
        enlarge y limits=0.04,
         ylabel style={at={(axis description cs:0,0.75)},anchor=east,
         rotate=-90},
         width=1.8in,
         height=1.3in
      },
    }
    \begin{tikzpicture}
    \begin{groupplot}[group style={group size=4 by 3,vertical sep=12pt}]
      \nextgroupplot[ylabel=(A)]
          \addplot+{2*x^2/(x^2+1)^2};

      \nextgroupplot[ytick={0.3}, yticklabels={2},ylabel=(B)]
          \addplot+{2*x^2/(x^2+1)^2 + 0.3};
          
      \nextgroupplot[ytick={-0.3}, yticklabels={$-2$},ylabel=(C)]
          \addplot+{-2*x^2/(x^2+1)^2 - 0.3};

      \nextgroupplot[ylabel=(D)]
          \addplot+[domain=-3.5:3.5]{x^2/(1+x^2)};

      \nextgroupplot[xtick={-2,2},ylabel=(E)]
          \addplot+[domain=-3.5:3.5]{1/5-1/(1+x^2)};

      \nextgroupplot[xtick={-2,2},ylabel=(F)]
          \addplot+[domain=-3.5:3.5]{-1/5+1/(1+x^2)};

      \nextgroupplot[ylabel=(G)]
          \addplot+[domain=-3.5:3.5]{atan(x)/90 - 3*x/(5+15*x^2)};

      \nextgroupplot[ytick={0.3}, yticklabels={2},ylabel=(H)]
          \addplot+[domain=-3.5:3.5]{atan(x)/90 - 3*x/(5+15*x^2)+0.3};

      \nextgroupplot[xtick={2},ylabel=(I)]
          \addplot+[domain=-3.5:3.5]{atan(x-2)/90 - 3*(x-2)/(5+15*(x-2)^2)};

      \nextgroupplot[ylabel=(J)]
          \addplot+{0.45};

      \nextgroupplot[ylabel=(K)]
          \addplot+{-0.45};    
    \end{groupplot}
    \end{tikzpicture}

\pspace{\newpage}

% Practice for Final
% \item ({Problem Set 23}, {{\#5}})

%  (close to 21.3.14)
% \begin{problem}
%     A sewage gutter is to be constructed from a rectangular piece of sheet
%     metal 6 feet wide by folding up a 2-foot strip on each side by equal
%     angles.  Let $\theta$ be the angle between the sides and the vertical,
%     as shown in the figure below.  What should $\theta$ be so that the
%     gutter will carry the maximum amount of water?  (Be sure to justify
%     that your value of $\theta$ actually maximizes the amount of water
%     carried.)

%     \begin{tikzpicture}[x=2cm, y=2cm]
%       \pgfmathsetmacro{\angle}{40}

%       \draw ({90+\angle}: 1) -- node[anchor=north east]{2} (0, 0) -- (1,
%       0) -- node[anchor=north west]{2} +({90-\angle}:1); %
%       \draw[dashed] (0, 0) -- (0, 1/3); %
%       \draw[dashed] (1, 0) -- (1, 1/3); %
%       \draw (0, 8pt) arc[radius=8pt, start angle=90, end
%       angle={90+\angle}]; %
%       \draw (1, 8pt) arc[radius=8pt, start angle=90, end
%       angle={90-\angle}]; %
%       \draw ({90+\angle/2}:12pt) node{$\theta$};
%       \draw ($(1, 0) + ({90-\angle/2}:12pt)$) node{$\theta$};
%     \end{tikzpicture}
%   \end{problem}

%   \begin{solution}
%     \ideas{We start by identifying our goal.}

%     \highlight{\textbf{Goal:} Maximize the amount of water the gutter can
%       carry, which is accomplished by maximizing the cross-sectional area
%       of the gutter.}

%     \ideas{Next, we write a formula for the thing we're trying to maximize
%       (cross-sectional area).  To do this, we'll need to define some
%       variables.}

%     \highlight{\textbf{Define variables:} Let $A$ be the cross-sectional
%       area of the gutter and $x, y$ be as labeled in the picture below.}

%     \begin{center}
%       \begin{tikzpicture}[x=2cm, y=2cm]
%         \pgfmathsetmacro{\angle}{40}

%         \draw ({90+\angle}: 1) -- node[anchor=north east]{2} (0, 0) --
%         node[below]{2} (1, 0) -- node[anchor=north west]{2}
%         +({90-\angle}:1) -- cycle; %
%         \draw[dashed] (0, 0) -- node[right]{$y$} (0, {cos(\angle)}); %
%         \draw[dashed] (1, 0) -- (1, {cos(\angle)}); %
%         \draw[decorate, decoration={brace}, yshift=2pt] ({90+\angle}:1) --
%         node[above, yshift=2pt]{$x$} (0, {cos(\angle)}); %
%         \draw (0, 8pt) arc[radius=8pt, start angle=90, end
%         angle={90+\angle}]; %
%         \draw ({90+\angle/2}:12pt) node{$\theta$};
%       \end{tikzpicture}
%     \end{center}
%     We can see that the area of the trapezoid can be broken up into two
%     triangles (of width $x$ and height $y$) and a rectangle (of width $2$
%     and height $y$).

%     \highlight{\textbf{Formula for thing we're trying to maximize:} $A = 2
%       \left( \frac{1}{2} x y \right) + 2 y = xy + 2y$}

%     \ideas{Next, we want to express $A$ as a function of just one
%       variable.  To do this, we must get $x$ and $y$ in terms of a single
%       variable.}  Since we're ultimately looking for $\theta$, let's
%     express $x$ and $y$ in terms of $\theta$.  Looking at the right
%     triangles in our picture, we see $\cos \theta = \frac{y}{2}$ and $\sin
%     \theta = \frac{x}{2}$, so $y = 2 \cos \theta$ and $x = 2 \sin \theta$.
%     Plugging this into our formula for $A$, we get $A$ in terms of just
%     $\theta$: $A(\theta) = 4 \sin \theta \cos \theta + 4 \cos \theta$.

%     \ideas{Now, maximize!}  In our scenario, the relevant values of
%     $\theta$ are between $0$ and $\frac{\pi}{2}$, so we need to maximize
%     $A(\theta)$ on $\left[ 0, \frac{\pi}{2} \right]$.  We first find the
%     derivative of $A(\theta)$:
%     \begin{align*}
%       A'(\theta) = 4 \cos^2 \theta - 4 \sin^2 \theta - 4 \sin \theta
%     \end{align*}
%     This is never undefined, and it's 0 when 
%     \begin{align*}
%       4 \cos^2 \theta - 4 \sin^2 \theta - 4 \sin \theta &= 0 \\
%       \shortintertext{Dividing through by 4,}
%       \cos^2 \theta - \sin^2 \theta - \sin \theta &= 0 \\
%       \intertext{By the Pythagorean Identity, $\sin^2 \theta + \cos^2
%         \theta = 1$, so $\cos^2 \theta = 1 - \sin^2 \theta$:}
%       1 - 2 \sin^2 \theta - \sin \theta &= 0 \\
%       \intertext{Now that the entire equation is in terms of $\sin
%         \theta$, let's substitute $u = \sin \theta$:}
%       1 - 2 u^2 - u &= 0 \\
%       - (2u^2 + u - 1) &= 0 \\
%       - (2u - 1)(u + 1) &= 0 \\
%       u & = \frac{1}{2}, -1 \\
%       \sin \theta &= \frac{1}{2}, -1 
%     \end{align*}
%     On $\left[ 0, \frac{\pi}{2} \right]$, the only solution is $\theta =
%     \frac{\pi}{6}$.  

%     \ideas{But wait, we're not done yet!  We still need to determine
%       whether this is the absolute maximum.}  Since our domain is a closed
%     interval, we can plug $\theta = \frac{\pi}{6}$, along with the
%     endpoints of the interval, into $A(\theta) = 4 \sin \theta \cos \theta
%     + 4 \cos \theta$ to see which gives us the highest value:
%     \begin{align*}
%       A(0) &= 4 \\
%       A \left( \frac{\pi}{6} \right) &= 4 \left( \frac{1}{2} \right)
%       \left( \frac{\sqrt{3}}{2} \right) + 4 \left( \frac{\sqrt{3}}{2}
%       \right) = \sqrt{3} + 2 \sqrt{3} = 3 \sqrt{3} \\
%       A \left( \frac{\pi}{2} \right) &= 0 
%     \end{align*}
%     So, the absolute maximum occurs when $\theta =
%     \boxed{\frac{\pi}{6}}$.     
%   \end{solution}

\begin{enumerate}[resume]

\item 

\begin{grading}
5 points: 1/4
\end{grading}

{\bf Reflect Back.} 

  \begin{enumerate}
  \item
\begin{problem}[\vfill]
      The area function 
      $A(x)=\displaystyle\int_a^x f(t)\,dt$ always has at least one $x$-intercept:
      at what $x$-value?
    \end{problem}

    \begin{solution}
      $A(x)$ has an $x$-intercept at $x = \boxed{a}$ because $
      A(a) = \displaystyle \int_a^a f(t)\,dt = 0$.
    \end{solution}

  \item
Last semester, you learned
    that the behavior of a function $F(x)$ is very much related to the
    behavior of its derivative.  Fill in the table below.  (If you know
    $F(x)$ has the property in the first column, what does that tell you
    about $F'$?)

    \begin{center}
      \renewcommand{\arraystretch}{2}
      \begin{tabular}{| p{1.25in} | p{1.25in} |}
        \hline
        \multicolumn{1}{|c|}{$F(x)$} & \multicolumn{1}{c|}{$F'(x)$} \\
        \hline\hline
        increasing & \solonly{\textbf{positive}}\\
        \hline
        decreasing & \solonly{\textbf{negative}} \\
        \hline
        concave up & \solonly{\textbf{increasing}} \\
        \hline
        concave down & \solonly{\textbf{decreasing}} \\
        \hline
      \end{tabular}      
    \end{center}

    \begin{problem}
      If $F(x) = A(x)$, what is $F'(x)$?
    \end{problem}

    \begin{solution}
      $F'(x)=f(x)$.
    \end{solution}
    
      \emph{Your table now tells you how the properties of $A$ relate
        to properties of $f$.}

        \pspace{\vfill}
  \end{enumerate}
\end{enumerate}

\psreflection{
For next class, read \S 23.3, 24.1.
}
\end{document}